\documentclass[11pt]{article}
\usepackage[margin=1in]{geometry}
\usepackage{xcolor}
\usepackage{mdframed}
\usepackage{enumitem}
\usepackage{hyperref}
\usepackage{booktabs}
\usepackage{parskip}

% Reviewer comment box
\newmdenv[
  backgroundcolor=gray!10,
  linecolor=gray!50,
  linewidth=0.8pt,
  innerleftmargin=10pt,
  innerrightmargin=10pt,
  innertopmargin=6pt,
  innerbottommargin=6pt,
  skipabove=6pt,
  skipbelow=6pt
]{reviewbox}

\newcommand{\reviewer}[1]{\vspace{1em}\noindent\textbf{\large #1}\vspace{0.4em}\hrule\vspace{0.6em}}
\newcommand{\comment}[2]{\vspace{0.8em}\noindent\textbf{Comment #1.}\begin{reviewbox}\textit{#2}\end{reviewbox}}
\newcommand{\response}{\noindent\textbf{Response:}~}
\newcommand{\revised}[1]{\textcolor{blue}{``#1''}}
\newcommand{\pending}[1]{\textcolor{red!70!black}{\textbf{[Pending — simulation required]:} #1}}

\title{\textbf{Response to Reviewer Comments}\\[0.5em]
\large COMMAG-26-00357\\
\normalsize \textit{End-to-End Performance of Video Streaming With MPEG-DASH Over Satellite 5G IAB Networks}}
\author{Muhammad Adeel Zahid, Ekram Hossain, Peng Hu}
\date{\today}

\begin{document}
\maketitle

\noindent We thank the Editor and all four reviewers for their thorough reading of the manuscript and for their constructive comments. The feedback has led to significant improvements to the paper. We address each comment below. Changes made to the manuscript are indicated in blue quoted excerpts; items that require new simulation runs are marked in red and will be incorporated in the final revised manuscript before resubmission.

\bigskip
\noindent\textbf{Summary of major changes:}
\begin{itemize}[noitemsep]
  \item Added orbital geometry analysis justifying why inter-satellite handover does not occur within the 60~s simulation window, and clarified that the satellite node already employs a waypoint mobility model.
  \item Replaced FDASH with BOLA (Buffer Occupancy-Based Lyapunov Algorithm) as the adaptive bitrate algorithm. All simulation results will be regenerated with BOLA and reported in the revised manuscript.
  \item Added a transport-layer analysis subsection explaining the interaction between congestion window dynamics and the ABR controller, including an explanation of the QUIC-NewReno throughput/bitrate anomaly.
  \item Expanded the Introduction to include 3GPP Release~17/18 NTN context, explicitly cite the gap between IAB and LEO satellite backhaul in current standards, and contrast this work against prior single-hop QUIC-over-satellite studies.
  \item Updated the bibliography with five new references: Leyva-Mayorga et al.\ (2020), Rinaldi et al.\ (2020), 3GPP TR~38.821, Spiteri et al.\ (2016) (BOLA), and Lai et al.\ (SIGCOMM 2025, LeoCC).
  \item Renamed stall metrics throughout to ``total stall duration'' and ``stall count'' for consistency with standard DASH QoE terminology.
  \item Added first-use definitions for all flagged acronyms: DASH, fDASH, HOL, gNB, CQI, RBG.
\end{itemize}

% ============================================================
\reviewer{Editor Comments (Dr.\ Lili Wei)}
% ============================================================

\comment{E1}{The mobility and handovers challenge of LEO-based backhaul in the full 60-second simulation is ignored, which is not a practical assumption.}

\response We respectfully clarify that the satellite node in our simulation already employs a \texttt{WaypointMobilityModel} that tracks its orbital movement at 7.8~km/s, and the 3GPP TR~38.811 NTN channel model computes propagation from this time-varying position, accounting for Doppler shift and changing path length throughout the simulation. The concern about handover is addressed by the following geometric analysis.

At an orbital altitude of 550~km, the satellite begins the simulation at the zenith above the simulation area (elevation angle $= 90°$). Over the 60~s simulation window it traverses approximately 468~km of ground track, reducing the elevation angle to approximately $50°$ — well above the minimum operational threshold of $5$--$10°$. Inter-satellite handover therefore does not occur within this window. The 60~s duration is representative of the highest-quality portion of a LEO satellite pass, where link geometry is strongest and transport behavior is least confounded by handover disruptions. We have added a paragraph to the Performance Evaluation setup making this argument explicit and providing the geometric calculation.

Regarding UE mobility: the scenario is a rural line-of-sight environment with no obstacles or reflectors, so pedestrian or vehicular UE movement at a few m/s does not meaningfully alter the access link geometry. This is noted in the revised manuscript.

The revised manuscript now reads:

\revised{The satellite node employs a waypoint mobility model that tracks its orbital movement at 7.8~km/s throughout the simulation, and the 3GPP TR 38.811 NTN channel model computes propagation characteristics from this time-varying position, accounting for Doppler shift and the changing path length. At an altitude of 550~km, the satellite is placed at the zenith above the simulation area (elevation angle $= 90°$) at $t = 0$. Over the 60~s simulation window it traverses approximately 468~km of ground track, reducing the elevation angle to approximately $50°$ — well above the minimum operational threshold of $5$--$10°$. Inter-satellite backhaul handover therefore does not occur within this window, and the scenario is representative of the strongest portion of a LEO satellite pass, where link quality is highest.}

We have also updated the Conclusion to include a more concrete future-work item describing how to extend the framework to simulate a full satellite pass (approximately 4--7~minutes above a 10° elevation threshold) to capture handover boundary effects.

\comment{E2}{It is unclear why FDASH was selected over other ABR schemes.}

\response We agree that the original manuscript did not adequately justify the choice of FDASH. Rather than providing a post-hoc defense, we have decided to replace FDASH with the BOLA (Buffer Occupancy-Based Lyapunov Algorithm) of Spiteri et al.\ (INFOCOM 2016), which is a well-established, widely-cited ABR scheme with a clear optimality formulation. We are implementing BOLA in the simulation framework and will re-run all five protocol configurations (TCP-BBR, TCP-CUBIC, TCP-NewReno, QUIC-BBR, QUIC-NewReno) with 30 seeds each. The revised manuscript will replace all figures and tables with BOLA-based results.

\pending{BOLA implementation in \texttt{src/dash/model/dash-client.cc}; full re-run of all configurations; replacement of Figures 2--3 and Tables I--II.}

\comment{E3}{Insufficient discussion of transport-layer behavior, including congestion window dynamics and RTT variation.}

\response The simulation already logs congestion window (cwnd) evolution and RTT time series for all configurations via trace callbacks connected to the QUIC and TCP socket layers. We have added a transport-layer analysis paragraph to the Discussion section that explains how cwnd and RTT dynamics connect to the observed application-level outcomes. In particular, we explain the QUIC-NewReno anomaly (high playback bitrate despite low throughput and high stall duration) through the interaction between its loss-based window reductions and the ABR controller's bitrate selection logic. Figures showing cwnd and RTT traces over time will be added once the BOLA-based re-runs are complete.

\pending{cwnd and RTT figures extracted from re-run simulation logs; added as a new figure in the Performance Evaluation section.}

% ============================================================
\reviewer{Reviewer 1}
% ============================================================

\comment{R1.1}{My main concern is the novelty of this work. The performance of QUIC over satellite networks have already been tested before, and hence, a comparison of congestion control algorithms is an incremental improvement.}

\response We agree that QUIC over satellite links has been studied (e.g., Alshagri et al.\ 2023; Yang et al.\ 2018). However, those works evaluate QUIC over direct satellite links (single-hop point-to-point paths). The present work introduces a qualitatively different architecture: a \emph{two-hop} 5G IAB path in which a LEO satellite provides the IAB backhaul and a 28~GHz mmWave link connects the IAB relay node to user equipment. This two-hop structure creates a distinct RTT composition, a unique interference pattern from the multi-beam satellite operating alongside the mmWave access link, and a bottleneck location that shifts depending on channel conditions on each hop. These characteristics have not been studied before, and they produce transport behavior (e.g., the QUIC-NewReno bitrate anomaly explained in Comment E3) that is specific to this architecture and would not arise in single-hop satellite evaluations.

We have rewritten the Introduction to make this distinction explicit, and we now contrast this work directly against prior single-hop studies and against the recent LeoCC work (Lai et al., SIGCOMM 2025) which addresses congestion control for direct LEO paths.

The revised manuscript now reads:

\revised{Unlike prior evaluations of QUIC over direct satellite links, this work examines a two-hop architecture in which a LEO satellite provides the IAB backhaul and a 28~GHz mmWave link connects the IAB node to user equipment. This two-hop path creates a distinct RTT composition and interference structure not present in single-hop satellite scenarios, and its interaction with the MPEG-DASH adaptive bitrate (ABR) logic has not previously been characterized.}

\comment{R1.2}{UE Mobility: The current simulation assumes stationary users. Real-world 5G users are mobile, which would introduce more frequent handovers and signal fluctuations.}

\response Please see Response~E1. The rural line-of-sight scenario (3GPP TR~38.901, RMa channel model) is free of obstacles and reflectors, so UE movement at pedestrian or vehicular speeds does not introduce signal fluctuations through diffraction or multipath effects. Handovers between IAB nodes are not triggered within the 60~s window because all UEs remain within the coverage radius of the single IAB node. We have added a sentence in the Performance Evaluation section noting this.

\comment{R1.3}{Satellite Dynamics: The serving satellite is assumed to stay in coverage for the full 60s. Incorporating constellation-level dynamics (satellite-to-satellite handovers) would add significant complexity and realism.}

\response Please see Response~E1. The 60~s window corresponds to a satellite elevation angle range of $90°$ to approximately $50°$, which is within continuous coverage. We agree that constellation-level handover dynamics are important and have added this as a concrete future-work item in the Conclusion, describing the simulation extension needed (two satellites with waypoint mobility, handover trigger at a minimum elevation threshold, 4--7~min simulation duration).

\comment{R1.4}{Cross-Layer Optimization: There is an opportunity to use link-layer information as feedback to the transport layer about impending handovers to prevent congestion window collapses.}

\response We agree this is a valuable direction. We have updated the Conclusion's future-work bullet to specifically describe this: \revised{Incorporate cross-layer information, such as predicted satellite elevation angle and link quality indicators, to provide the transport layer with advance warning of impending handovers and prevent congestion window collapses.}

% ============================================================
\reviewer{Reviewer 2}
% ============================================================

\comment{R2.1}{The authors assume the serving satellite remains in coverage for the full 60-second simulation. As we know, LEO satellites move at approximately 7.5 km/s. In a 60-second window, a satellite would travel 450 km resulting in a link handover. This paper ignores mobility and handovers and fails to address the most critical challenge of LEO-based backhaul.}

\response We respectfully note that ground-track distance and link handover are not equivalent. Whether handover occurs depends on the satellite's elevation angle at the start of the window, not on how far it travels. In our simulation, the satellite starts at the zenith (elevation $= 90°$). After 60~s of movement at 7.8~km/s, the ground-track displacement is 468~km, which shifts the elevation angle to approximately $50°$. The satellite remains well above the minimum operational elevation threshold ($5$--$10°$). A handover would only occur if the satellite were already near the horizon at the start of the simulation window, which is not the case here.

Furthermore, the satellite node in our simulation already employs a \texttt{WaypointMobilityModel} at 7.8~km/s, and the 3GPP TR~38.811 channel model accounts for the resulting Doppler shift and path-length variation in every simulation step. The channel is therefore not static.

We have added a paragraph to the Performance Evaluation section containing the elevation angle calculation and clarifying the above reasoning. Please see the updated manuscript.

\comment{R2.2}{The simulated topology is somewhat simple, with only one IAB node. The reviewer suggested that the authors conduct experiments on a more realistic topology.}

\response The single-IAB-node topology is a deliberate and well-motivated design choice. As described in the Network Architecture section, adding more IAB nodes in this scenario introduces inter-beam interference from the multibeam satellite and degrades link quality, ultimately causing some users to lose connectivity entirely. This is a property specific to the multibeam satellite + IAB combination and is not present in single-hop satellite studies. The single-IAB topology isolates the effect of the transport protocol and congestion control algorithm, which is the focus of this paper. We have retained the explanation already in the manuscript and ensured it is clearly stated.

\comment{R2.3}{The authors need to update the bibliography. I suggest to cite more relevant works.}

\response We have added five new references: Leyva-Mayorga et al.\ (IEEE Access, 2020) on LEO constellations for 5G; Rinaldi et al.\ (IEEE Access, 2020) on NTN in 5G and beyond; 3GPP TR~38.821 (NR NTN solutions, Release~16/17); Spiteri et al.\ (INFOCOM 2016) on BOLA; and Lai et al.\ (SIGCOMM 2025) on LeoCC congestion control for LEO satellite dynamics.

% ============================================================
\reviewer{Reviewer 3}
% ============================================================

\comment{R3.1}{In the introduction section, why do they claim that IAB architectures with LEO satellite backhaul are not currently addressed? Why is there no analysis of what is currently specified regarding Non-Terrestrial Networks (NTN) described in 3GPP Releases 17 and 18?}

\response We have added a paragraph to the Introduction addressing this directly. The paragraph now explains that 3GPP Releases~17 and~18 introduced NTN specifications (TR~38.821, TR~23.700-27) for direct satellite-to-device connectivity, but that the specific combination of 5G IAB relay operation with a LEO satellite backhaul link is not addressed in any current standard. The revised text reads:

\revised{5G IAB was first standardized in 3GPP Release~16, enabling the 5G New Radio (NR) spectrum to serve both UE access links and wireless backhaul among next-generation NodeBs (gNBs). Releases~17 and~18 introduced NTN support for direct satellite-to-device connectivity, and recent work has explored congestion control adaptation to LEO dynamics in general satellite paths. However, the combination of 5G IAB with LEO satellite backhaul is not currently addressed by 3GPP standards, and no prior study has jointly evaluated transport protocol behavior and MPEG-DASH QoE in this two-hop integrated architecture.}

\comment{R3.2}{The use of FDASH is defined. Why was the use of a fuzzy logic algorithm preferred without comparing it with standard ABR algorithms such as BOLA? How can they ensure that the high number of interruptions is due to the transport protocols and not a suboptimal reaction of the FDASH controller itself?}

\response This is an important concern. We have replaced FDASH with BOLA as the ABR algorithm for the revised manuscript. BOLA has a well-defined optimality criterion (it approximately maximizes a Lyapunov utility function of bitrate and buffer occupancy) and is widely used as a reference ABR scheme in the literature. Replacing FDASH with BOLA removes the ambiguity about whether observed stalls are caused by the transport protocol or by an overly aggressive or overly conservative ABR controller. All simulation results will be regenerated with BOLA.

\pending{BOLA implementation and full re-run of all five protocol configurations.}

\comment{R3.3}{It is stated that the UE remains stationary and that the satellite assumes continuous coverage throughout the simulation. LEO satellites are characterized by their high orbital velocity and frequent handovers. Why was it decided to exclude mobility and line-of-sight loss events from the simulation?}

\response Please see Responses E1 and R2.1. The satellite is not static — it moves via \texttt{WaypointMobilityModel} at 7.8~km/s, and the channel model accounts for the resulting Doppler and geometry changes. The elevation angle at the end of the 60~s window is approximately $50°$, which is within continuous coverage. UE mobility is negligible in a rural LOS scenario. These points are now stated explicitly in the revised manuscript.

\comment{R3.4}{In Table II it is defined that the playback interval is 60 seconds. To evaluate congestion control algorithms, congestion windows require much longer time frames to reach convergence. Do these results not merely reflect the initial or transient phase of the connections instead of their sustained steady-state behavior?}

\response This is a valid observation for TCP in a standard terrestrial setting. However, in a LEO satellite environment with a one-way propagation delay of approximately 2~ms (for a 550~km altitude link) and an end-to-end RTT determined by the two-hop path, the cwnd evolution is faster than in GEO satellite scenarios where RTTs exceed 500~ms. We have addressed this in the Discussion section by adding an analysis of cwnd and RTT dynamics, explaining how the transient and steady-state phases manifest for each protocol configuration in this specific network environment. This analysis shows that the 60~s window is sufficient to observe meaningful differentiation between configurations, though we acknowledge that longer simulations would be needed to study sustained steady-state behavior at lower elevation angles or during handover events.

The concern also motivates our future-work item on extending the simulation to cover a full satellite pass ($\sim$4--7~min), which would provide both steady-state behavior and handover disruption data.

\comment{R3.5}{Formatting: In Figures 2 and 3, the names of the algorithms appear without a hyphen (QUIC BBR) while in most of the text a hyphen is used (QUIC-BBR).}

\response We will correct the figure labels to use ``QUIC-BBR'' and ``QUIC-NewReno'' consistently when regenerating the figures for the BOLA-based re-run.

\pending{Figure label correction when plots are regenerated.}

\comment{R3.6}{Several acronyms need to be defined: DASH, fDASH, HOL, gNB, CQI, RB, RBG, and others.}

\response We have added first-use definitions for all acronyms that appear in the manuscript body: MPEG-DASH (Dynamic Adaptive Streaming over HTTP) is defined in the Introduction; FDASH (Fuzzy DASH) is defined in the Application Layer subsection; HOL (head-of-line) blocking is defined in the Background section; gNB (next-generation NodeB) is defined in the Introduction; CQI (Channel Quality Indicator) and RBG (Resource Block Group) are defined in Table~I. The acronyms DU, MT, NAS, RRC, PDCP, and RLC listed by the reviewer do not appear in the article body and therefore do not require definition.

% ============================================================
\reviewer{Reviewer 4}
% ============================================================

\comment{R4.1}{Many recent studies have already investigated transport protocols and video streaming over LEO networks. The novelty of this paper is not yet sufficiently clear. The unique contribution should be its attempt to identify previously unexplored challenges in ``5G-IAB-integrated'' LEO networks.}

\response Please see Response R1.1. We have rewritten the Introduction to make the two-hop IAB-specific contribution explicit. The revised text emphasizes three aspects unique to the IAB-LEO architecture that distinguish this work from prior single-hop satellite studies: (1) the 28~GHz mmWave access hop introduces a second delay and loss source with different statistical properties than the satellite hop; (2) the IAB relay creates multi-beam interference between the satellite backhaul and the terrestrial access links; and (3) the resulting RTT composition and bottleneck location produce transport behaviors (documented in the Discussion) that do not arise in direct satellite links.

\comment{R4.2}{There is insufficient discussion of transport-layer behavior, including congestion window dynamics and RTT variation. I strongly recommend reorganizing the performance evaluation section to better connect the observed application-level results with the underlying transport-layer mechanisms.}

\response Please see Response E3. We have added a transport-layer analysis paragraph to the Discussion. The revised manuscript explains that TCP-CUBIC and TCP-NewReno build large queues at the IAB backhaul bottleneck, resulting in high RTT (approximately 289~ms) that makes them unsuitable for delay-sensitive streaming despite adequate throughput. QUIC-BBR's model-based cwnd management avoids queue buildup, yielding substantially lower RTT (12.65~ms) and the best balance of stall duration and latency. The QUIC-NewReno anomaly (high bitrate, high stalls) is explained through the interaction between aggressive bitrate selection by the ABR controller under low RTT conditions and the periodic cwnd collapses from loss-based congestion control.

Figures showing cwnd and RTT time series will be added in the revised manuscript once the BOLA-based simulations are complete.

\pending{cwnd and RTT time-series figures for all five configurations.}

\comment{R4.3}{The rationale for adopting FDASH as the ABR method is unclear.}

\response As noted in Response~E2, we are replacing FDASH with BOLA. This eliminates the need to justify FDASH and provides a well-established, theoretically grounded ABR baseline.

\comment{R4.4}{Figure 3 evaluates playback duration as a performance metric. It is not clear why stalling time was not used instead, as it is a more common QoE-related metric in DASH studies.}

\response We agree. We report stalling time (referred to as ``total stall duration'' in Table~II) in addition to playback duration, but the connection between the two was not made clear. We have renamed the metric throughout the manuscript as ``total stall duration'' (the cumulative time the playback is paused), added a note that this is the standard stall metric used in DASH QoE studies, and added an explanation of why playback duration is also reported (it captures whether the session completes at all — a binary QoE outcome relevant to satellite disruption scenarios). The table caption, column headers, and all prose references have been updated accordingly.

\comment{R4.5}{The current manuscript does not sufficiently explain why QUIC/NewReno shows better performance in terms of throughput. This discrepancy with prior understanding should be discussed more carefully.}

\response This is addressed in the new transport-layer analysis paragraph added to the Discussion (see Response~E3). QUIC-NewReno achieves the highest median playback bitrate (3.28~Mbps) not because of high sustained throughput, but because its conservative cwnd behavior (frequent halving on loss events) keeps queue occupancy low, producing the lowest RTT (7.48~ms) among all tested configurations. The FDASH — and in the revised manuscript, BOLA — controller interprets the low RTT and fast per-segment delivery as a signal of available bandwidth and selects a high bitrate. However, the small cwnd limits the sustained rate, so the buffer drains faster than it fills during high-bitrate phases, causing frequent and prolonged stalls (34.33~s median total stall duration). This ``high bitrate, high stalls'' pattern is specific to the IAB two-hop bottleneck and is explained in detail in the Discussion.

\bigskip
\hrule
\bigskip

\noindent\textbf{Note on pending items:} The items marked \textcolor{red!70!black}{\textbf{[Pending — simulation required]}} above (BOLA implementation and re-run, cwnd/RTT figures, figure label corrections) will be completed and incorporated before the revised manuscript is submitted. The writing-level changes described in this letter are already reflected in the current draft.

\end{document}
